\section{Discussion: From Coalition Effects to Coalition Policy}\label{sec:discussion}

\subsection{The organizational question resolved by the model}

The paper begins with a routing-control question rather than a concentration
question. AV companies face a common physical congestion technology but can
use different management objectives for their controlled vehicles. Coalition
governance determines which of those objectives are optimized jointly and
which cross-member effects are internalized. Because each member retains its
accepted OD demand, an organizational change first alters an OD-feasible route
response. The road network maps that response into edge flows, and HDVs then
absorb the component they can reproduce through Wardrop reallocation. Only the
remaining activated response can change aggregate congestion.

This mechanism explains why coalition degree lacks a universal ordering. A
larger coalition need not generate a larger or more socially aligned activated
response. Its cross-member terms can rotate the response toward a different
set of congested edges, and HDV substitution can remove some directions but
not others. The relevant policy object is therefore the system-cost projection
of the screened response. Proposition~\ref{prop:merge-regimes} makes this
statement operational: a proposed increase in coupling improves TSTT when its
merge-direction score is negative, is neutral when the score is zero, and is
harmful when the score is positive.

\subsection{When centralization and partial independence are preferred}

The canonical theorem turns the general directional criterion into a regime
map. When interior HDVs can offset the coalition response, they pin the
congested route at the Wardrop flow and organization is locally neutral. Once
that screening support disappears, the fixed component $H$ and effective
coalition response $\Sigma$ determine system performance. If
$H\geq x^{UE}/2$, TSTT is nondecreasing in $\Sigma$, so a centralizing move is
desirable when it reduces effective response. If $H<x^{UE}/2$, system cost is
minimized at the interior target
$\Sigma^*=1-2H/x^{UE}$. A partial coalition can then dominate both complete
centralization and complete fragmentation because it produces a response
closer to the target.

The theorem does not assert that a merger always reduces $\Sigma$. That
relationship depends on the declared objective family and governance matrix.
It instead states a sharper condition: once the response induced by an
organization is known or estimated, the preferred direction follows from its
position relative to the active regime. On a general network, the scalar
target becomes the merge-direction score, which must be recomputed after a
support change.

\subsection{What the network evidence adds}

The Sioux Falls audit scan is consistent with the regime logic. At
$\alpha=0.5$, the grand coalition minimizes TSTT for both tested positive
governance intensities. At $\alpha=0.9$, the optimal set retains independent
decision centers: symmetry-equivalent three--one partitions at $\gamma=0.5$
and a two--one--one partition at $\gamma=1$. The two-network validation shows
the same qualitative transition. All tested Braess and grid states with
$\alpha\leq0.7$ select the grand coalition, whereas high penetration and
sufficiently strong governance produce partial optima on both networks. The
transition occurs at different values of $\gamma$, so topology changes the
boundary without eliminating the regime.

These results improve upon a list of partition rankings in two ways. First,
the analytic theorem explains why an interior organization can be optimal.
Second, policy regret asks whether a decision rule can recover the organization
without treating the full partition lattice as the prescription. One-step
merging is exact in all four Sioux Falls states and $29$ of $32$ cross-network
states, with only $0.332$ points of mean cross-network regret. Its three Braess
failures are informative: local improvement can stop before a jointly valuable
sequence of merges becomes visible. Two-step lookahead crosses these local
traps and reaches an audit optimum in all $36$ states. This observed exactness
does not establish universal optimality, particularly because the four-company
two-step neighborhood covers much of the small partition lattice.

\subsection{Governance burden changes the recommendation}

Minimizing TSTT alone treats coordination as costless. The burden model in
\eqref{eq:governance-burden} separates the network benefit from the resources
needed to maintain cross-company governance. Each arrangement then defines an
affine loss in $\lambda$: its congestion performance is the intercept and its
coordination burden is the slope. The lower envelope yields exact switching
thresholds. In the Sioux Falls example at $\alpha=0.5,\gamma=1$, the selected
organization moves from grand coordination to two pairs, then to one pair, and
finally to singletons as the governance-cost weight rises.

This result gives partial coalitions a substantive interpretation. They are not
arbitrary intermediate points or evidence that moderation is always best.
They appear when the marginal congestion benefit of an additional cross-member
link no longer covers its governance burden. Alternative cost functions could
change the thresholds or remove an intermediate interval; the lower-envelope
method remains valid once that burden is specified.

\subsection{Why HHI cannot implement the policy}

The response-based policy needs variables that a passive ownership cross
section does not reveal. The controlled composition experiment holds company
and coalition HHI, masses, coalition sizes, OD obligations, and route access
fixed, yet changing which objective types are grouped alters recovery and can
reverse the ordering along the governance path. Across the full Sioux Falls
audit, the coalition-HHI/recovery Spearman correlation falls to $0.236$ and the
pairwise inversion rate reaches $40.74\%$ at high penetration and full
governance. At fixed company HHI, changing spatial OD portfolios creates even
larger variation because it changes the network directions controlled by each
company.

HHI remains useful for describing the distribution of fleet mass. It cannot
select a coalition organization because it omits objective curvature,
cross-member governance, OD authority, route access, active support, and HDV
screening. The policy should therefore treat HHI as a screening covariate and
estimate or simulate the activated response directly.

\subsection{Observability and implementation}

Implementation requires four classes of inputs: accepted company-level OD
obligations, available routes or dispatch authority, calibrated physical delay
functions, and a defensible approximation of each company's marginal
management objective. Under a stable support, OD-preserving perturbations can
identify screened aggregate-flow responses and provide the ingredients for the
merge score. When supports change, the manager can use the certified
continuation solver to evaluate the local one- or two-step merge neighborhood.
This creates an auditable workflow: diagnose the response regime, evaluate
candidate local merges, report regret against richer audit scans when feasible,
and include governance burden before recommending an organization.

The recommendation applies to an operator, platform, or regulator that can
compare specified governance arrangements. It does not imply that every member
would accept the system-optimal continuation design. Private implementation
requires a payoff-allocation or transfer rule, contracting costs, deviation
protocols, and a partition-formation concept.

\subsection{Scope and limitations}

The exact response and regime results require affine physical costs, fixed
active supports, nonredundant OD tangents, and positive curvature on
coalition-feasible directions. The BPR experiments provide certified numerical
continuation equilibria, not a global nonlinear extension of the affine
theorem. The governance burden is a transparent sensitivity primitive rather
than an externally calibrated institutional cost. Objective types are also
controlled primitives; estimating them from operator dispatch data remains an
important empirical task.

The routing stage conditions on an externally realized batch of accepted
requests. Pricing, passenger choice, endogenous demand, empty-vehicle
repositioning, service quality, capacity and headway effects, private profit,
and endogenous coalition stability remain outside the model. The complete
partition scans cover four-company benchmark menus rather than a population of
markets, and the two-step policy has no general optimality guarantee. These
boundaries leave a clear next step: combine empirically identified objectives
with scalable local search and a partition-function formation game, while
retaining the continuation-equilibrium mechanism established here.
